\section{早明的朝廷、法律与文化连结}

本节以朝廷程序、法律语言、家庭与性别、城市文化和海外联系为问题线索。正史、实录、律例与奏议通常更容易保存官员、男性精英和国家分类；女性、家庭成员、工匠、城市居民与海外中介则常仅在墓志、契约、地方志、碑刻、航海记录或案件中留下片段。材料本身的偏向，应当成为叙事的一部分。

\subsection{宫廷程序与法律语言}

\paragraph{1375（洪武八年）：明开始发行新纸币“大明宝钞”} 此事把权力、法律或文化关系推进到可辨认的节点。账本的证据限度是编年候选的年序可由官修与后出材料交叉；具体月日、参与者、战果或数字必须回查所列材料，不能将单一叙事当作定论。可资核查的材料为《太祖实录》洪武八年条；《明史·本纪》及相关列传；诏令、奏疏或会典版本。它们能够说明制度如何表述自身、某种命令或交涉如何被记录，却未必能直接复原家庭内部关系、性别分工、城市日常或海外行动者的意图。事件发生的条件涉及继承、官僚协调或皇权运作的压力使问题进入朝廷程序；其后留下的制度与社会后果为它影响后续人事与政策议程；动机和派系标签不能由单一叙事裁定。译写候选以编年材料定位；实录记载反映朝廷选择，崇祯期须另以奏疏、地方及敌对政权材料交叉。

\paragraph{1375（洪武八年）：由于费用和浪费，洪武皇帝停止了凤阳的建设；建设计划转移至南京} 此事把权力、法律或文化关系推进到可辨认的节点。账本的证据限度是编年候选的年序可由官修与后出材料交叉；具体月日、参与者、战果或数字必须回查所列材料，不能将单一叙事当作定论。可资核查的材料为《太祖实录》洪武八年条；《明史·本纪》及相关列传；诏令、奏疏或会典版本。它们能够说明制度如何表述自身、某种命令或交涉如何被记录，却未必能直接复原家庭内部关系、性别分工、城市日常或海外行动者的意图。事件发生的条件涉及继承、官僚协调或皇权运作的压力使问题进入朝廷程序；其后留下的制度与社会后果为它影响后续人事与政策议程；动机和派系标签不能由单一叙事裁定。译写候选以编年材料定位；实录记载反映朝廷选择，崇祯期须另以奏疏、地方及敌对政权材料交叉。

\paragraph{1375（洪武八年）：皇室、功臣与中枢官僚的任免或案件材料构成本年权力重组线索，案狱叙事须避免照录罪名} 此事把权力、法律或文化关系推进到可辨认的节点。账本的证据限度是来源导向的年度桥接条目：本年材料可定位相关政务线索；具体月日、发文机关、地域和执行结果仍须逐项回查，不能以制度文本或奏报替代实际效果。可资核查的材料为《太祖实录》洪武八年条；《明史·本纪》及相关列传；诏令、奏疏或会典版本。它们能够说明制度如何表述自身、某种命令或交涉如何被记录，却未必能直接复原家庭内部关系、性别分工、城市日常或海外行动者的意图。事件发生的条件涉及继承、官僚协调或皇权运作的压力使问题进入朝廷程序；其后留下的制度与社会后果为它影响后续人事与政策议程；动机和派系标签不能由单一叙事裁定。桥接条目只标示可回查的年度问题与材料链；实录有中央选择性，崇祯期尤其需要奏疏、地方、朝鲜及满文材料交叉。

\paragraph{1376（洪武九年）三月：明军击败巴彦帖木儿} 此事把权力、法律或文化关系推进到可辨认的节点。账本的证据限度是编年候选的年序可由官修与后出材料交叉；具体月日、参与者、战果或数字必须回查所列材料，不能将单一叙事当作定论。可资核查的材料为《太祖实录》洪武九年条；《明史·兵志》及相关列传；边镇、地方志或对方材料。它们能够说明制度如何表述自身、某种命令或交涉如何被记录，却未必能直接复原家庭内部关系、性别分工、城市日常或海外行动者的意图。事件发生的条件涉及前期边防、地方武装或跨区域资源调动的压力推动了该行动；其后留下的制度与社会后果为它改变了后续调兵、补给或地方秩序的条件；战果和伤亡须分开核对。译写候选以编年材料定位；实录记载反映朝廷选择，崇祯期须另以奏疏、地方及敌对政权材料交叉。

\paragraph{1376（洪武九年）七月：明军再次击败巴彦帖木儿} 此事把权力、法律或文化关系推进到可辨认的节点。账本的证据限度是编年候选的年序可由官修与后出材料交叉；具体月日、参与者、战果或数字必须回查所列材料，不能将单一叙事当作定论。可资核查的材料为《太祖实录》洪武九年条；《明史·兵志》及相关列传；边镇、地方志或对方材料。它们能够说明制度如何表述自身、某种命令或交涉如何被记录，却未必能直接复原家庭内部关系、性别分工、城市日常或海外行动者的意图。事件发生的条件涉及前期边防、地方武装或跨区域资源调动的压力推动了该行动；其后留下的制度与社会后果为它改变了后续调兵、补给或地方秩序的条件；战果和伤亡须分开核对。译写候选以编年材料定位；实录记载反映朝廷选择，崇祯期须另以奏疏、地方及敌对政权材料交叉。

\paragraph{1376（洪武九年）十月22日：洪武皇帝宣布他将接受官员对其统治的直接批评} 此事把权力、法律或文化关系推进到可辨认的节点。账本的证据限度是编年候选的年序可由官修与后出材料交叉；具体月日、参与者、战果或数字必须回查所列材料，不能将单一叙事当作定论。可资核查的材料为《太祖实录》洪武九年条；《明史·兵志》及相关列传；边镇、地方志或对方材料。它们能够说明制度如何表述自身、某种命令或交涉如何被记录，却未必能直接复原家庭内部关系、性别分工、城市日常或海外行动者的意图。事件发生的条件涉及前期边防、地方武装或跨区域资源调动的压力推动了该行动；其后留下的制度与社会后果为它改变了后续调兵、补给或地方秩序的条件；战果和伤亡须分开核对。译写候选以编年材料定位；实录记载反映朝廷选择，崇祯期须另以奏疏、地方及敌对政权材料交叉。

\subsection{宗室、家庭与礼制}

\paragraph{1376（洪武九年）：叶伯驹因批评皇帝重刑被饿死狱中} 此事把权力、法律或文化关系推进到可辨认的节点。账本的证据限度是编年候选的年序可由官修与后出材料交叉；具体月日、参与者、战果或数字必须回查所列材料，不能将单一叙事当作定论。可资核查的材料为《太祖实录》洪武九年条；《明史·本纪》及相关列传；诏令、奏疏或会典版本。它们能够说明制度如何表述自身、某种命令或交涉如何被记录，却未必能直接复原家庭内部关系、性别分工、城市日常或海外行动者的意图。事件发生的条件涉及继承、官僚协调或皇权运作的压力使问题进入朝廷程序；其后留下的制度与社会后果为它影响后续人事与政策议程；动机和派系标签不能由单一叙事裁定。译写候选以编年材料定位；实录记载反映朝廷选择，崇祯期须另以奏疏、地方及敌对政权材料交叉。

\paragraph{1376（洪武九年）：洪武帝处决所有与“预印文书案”有关的官员} 此事把权力、法律或文化关系推进到可辨认的节点。账本的证据限度是编年候选的年序可由官修与后出材料交叉；具体月日、参与者、战果或数字必须回查所列材料，不能将单一叙事当作定论。可资核查的材料为《太祖实录》洪武九年条；《明史·本纪》及相关列传；诏令、奏疏或会典版本。它们能够说明制度如何表述自身、某种命令或交涉如何被记录，却未必能直接复原家庭内部关系、性别分工、城市日常或海外行动者的意图。事件发生的条件涉及继承、官僚协调或皇权运作的压力使问题进入朝廷程序；其后留下的制度与社会后果为它影响后续人事与政策议程；动机和派系标签不能由单一叙事裁定。译写候选以编年材料定位；实录记载反映朝廷选择，崇祯期须另以奏疏、地方及敌对政权材料交叉。

\paragraph{1377（洪武十年）五月：明军入侵青海} 此事把权力、法律或文化关系推进到可辨认的节点。账本的证据限度是编年候选的年序可由官修与后出材料交叉；具体月日、参与者、战果或数字必须回查所列材料，不能将单一叙事当作定论。可资核查的材料为《太祖实录》洪武十年条；《明史·兵志》及相关列传；边镇、地方志或对方材料。它们能够说明制度如何表述自身、某种命令或交涉如何被记录，却未必能直接复原家庭内部关系、性别分工、城市日常或海外行动者的意图。事件发生的条件涉及前期边防、地方武装或跨区域资源调动的压力推动了该行动；其后留下的制度与社会后果为它改变了后续调兵、补给或地方秩序的条件；战果和伤亡须分开核对。译写候选以编年材料定位；实录记载反映朝廷选择，崇祯期须另以奏疏、地方及敌对政权材料交叉。

\paragraph{1377（洪武十年）：南京宫殿竣工，定为“京师”} 此事把权力、法律或文化关系推进到可辨认的节点。账本的证据限度是编年候选的年序可由官修与后出材料交叉；具体月日、参与者、战果或数字必须回查所列材料，不能将单一叙事当作定论。可资核查的材料为《太祖实录》洪武十年条；《明史·本纪》及相关列传；诏令、奏疏或会典版本。它们能够说明制度如何表述自身、某种命令或交涉如何被记录，却未必能直接复原家庭内部关系、性别分工、城市日常或海外行动者的意图。事件发生的条件涉及继承、官僚协调或皇权运作的压力使问题进入朝廷程序；其后留下的制度与社会后果为它影响后续人事与政策议程；动机和派系标签不能由单一叙事裁定。译写候选以编年材料定位；实录记载反映朝廷选择，崇祯期须另以奏疏、地方及敌对政权材料交叉。

\paragraph{1377（洪武十年）：朝廷围绕卫所军户的编籍、屯田与调发整理本年军政文书，作为明初军事接收的可核查节点} 此事把权力、法律或文化关系推进到可辨认的节点。账本的证据限度是来源导向的年度桥接条目：本年材料可定位相关政务线索；具体月日、发文机关、地域和执行结果仍须逐项回查，不能以制度文本或奏报替代实际效果。可资核查的材料为《太祖实录》洪武十年条；《明史·兵志》及相关列传；边镇、地方志或对方材料。它们能够说明制度如何表述自身、某种命令或交涉如何被记录，却未必能直接复原家庭内部关系、性别分工、城市日常或海外行动者的意图。事件发生的条件涉及前期边防、地方武装或跨区域资源调动的压力推动了该行动；其后留下的制度与社会后果为它改变了后续调兵、补给或地方秩序的条件；战果和伤亡须分开核对。桥接条目只标示可回查的年度问题与材料链；实录有中央选择性，崇祯期尤其需要奏疏、地方、朝鲜及满文材料交叉。

\paragraph{1378（洪武十一年）：吴冕叛乱：锦族叛乱} 此事把权力、法律或文化关系推进到可辨认的节点。账本的证据限度是编年候选的年序可由官修与后出材料交叉；具体月日、参与者、战果或数字必须回查所列材料，不能将单一叙事当作定论。可资核查的材料为《太祖实录》洪武十一年条；《明史·兵志》及相关列传；边镇、地方志或对方材料。它们能够说明制度如何表述自身、某种命令或交涉如何被记录，却未必能直接复原家庭内部关系、性别分工、城市日常或海外行动者的意图。事件发生的条件涉及前期边防、地方武装或跨区域资源调动的压力推动了该行动；其后留下的制度与社会后果为它改变了后续调兵、补给或地方秩序的条件；战果和伤亡须分开核对。译写候选以编年材料定位；实录记载反映朝廷选择，崇祯期须另以奏疏、地方及敌对政权材料交叉。

\subsection{都市空间与文化工程}

\paragraph{1378（洪武十一年）：学校、科举、礼制或律令的颁行与讨论在本年材料中留下制度线索，规定与地方实践应分开处理} 此事把权力、法律或文化关系推进到可辨认的节点。账本的证据限度是来源导向的年度桥接条目：本年材料可定位相关政务线索；具体月日、发文机关、地域和执行结果仍须逐项回查，不能以制度文本或奏报替代实际效果。可资核查的材料为《太祖实录》洪武十一年条；《明史·选举志》或《艺文志》；文集、碑刻或书目材料。它们能够说明制度如何表述自身、某种命令或交涉如何被记录，却未必能直接复原家庭内部关系、性别分工、城市日常或海外行动者的意图。事件发生的条件涉及制度、教育或知识传播的需要推动了相关行动或文本形成；其后留下的制度与社会后果为它提供制度和传播史的锚点，不能据此推断社会整体接受。桥接条目只标示可回查的年度问题与材料链；实录有中央选择性，崇祯期尤其需要奏疏、地方、朝鲜及满文材料交叉。

\paragraph{1378（洪武十一年）：户部、布政司与里甲体系围绕户籍、田土和赋役的申报形成年度行政链，本条标记其可回查节点} 此事把权力、法律或文化关系推进到可辨认的节点。账本的证据限度是来源导向的年度桥接条目：本年材料可定位相关政务线索；具体月日、发文机关、地域和执行结果仍须逐项回查，不能以制度文本或奏报替代实际效果。可资核查的材料为《太祖实录》洪武十一年条；《明会典》相关条目；地方志、黄册/契约/河工材料。它们能够说明制度如何表述自身、某种命令或交涉如何被记录，却未必能直接复原家庭内部关系、性别分工、城市日常或海外行动者的意图。事件发生的条件涉及财政征解、交通工程或货币支付的压力使相关措施进入政务；其后留下的制度与社会后果为其法定安排不等于地方执行，后续须以地域材料追踪负担与成效。桥接条目只标示可回查的年度问题与材料链；实录有中央选择性，崇祯期尤其需要奏疏、地方、朝鲜及满文材料交叉。

\paragraph{1379（洪武十二年）二月：明军在甘肃击败藏人} 此事把权力、法律或文化关系推进到可辨认的节点。账本的证据限度是编年候选的年序可由官修与后出材料交叉；具体月日、参与者、战果或数字必须回查所列材料，不能将单一叙事当作定论。可资核查的材料为《太祖实录》洪武十二年条；《明史·兵志》及相关列传；边镇、地方志或对方材料。它们能够说明制度如何表述自身、某种命令或交涉如何被记录，却未必能直接复原家庭内部关系、性别分工、城市日常或海外行动者的意图。事件发生的条件涉及前期边防、地方武装或跨区域资源调动的压力推动了该行动；其后留下的制度与社会后果为它改变了后续调兵、补给或地方秩序的条件；战果和伤亡须分开核对。译写候选以编年材料定位；实录记载反映朝廷选择，崇祯期须另以奏疏、地方及敌对政权材料交叉。

\paragraph{1379（洪武十二年）：占城向南京致敬} 此事把权力、法律或文化关系推进到可辨认的节点。账本的证据限度是编年候选的年序可由官修与后出材料交叉；具体月日、参与者、战果或数字必须回查所列材料，不能将单一叙事当作定论。可资核查的材料为《太祖实录》洪武十二年条；《明史·外国传》；《朝鲜王朝实录》、琉球《历代宝案》或海外档案。它们能够说明制度如何表述自身、某种命令或交涉如何被记录，却未必能直接复原家庭内部关系、性别分工、城市日常或海外行动者的意图。事件发生的条件涉及朝贡名义、边境安全与港口贸易的利益使交涉持续发生；其后留下的制度与社会后果为它为后续册封、互市或海防安排留下文书与跨政权比较线索。译写候选以编年材料定位；实录记载反映朝廷选择，崇祯期须另以奏疏、地方及敌对政权材料交叉。

\paragraph{1379（洪武十二年）：北元诸部、东北和西南的边报、招抚与防御命令持续进入本年政务，须按具体部众和地域复核} 此事把权力、法律或文化关系推进到可辨认的节点。账本的证据限度是来源导向的年度桥接条目：本年材料可定位相关政务线索；具体月日、发文机关、地域和执行结果仍须逐项回查，不能以制度文本或奏报替代实际效果。可资核查的材料为《太祖实录》洪武十二年条；《明史·外国传》；《朝鲜王朝实录》、琉球《历代宝案》或海外档案。它们能够说明制度如何表述自身、某种命令或交涉如何被记录，却未必能直接复原家庭内部关系、性别分工、城市日常或海外行动者的意图。事件发生的条件涉及朝贡名义、边境安全与港口贸易的利益使交涉持续发生；其后留下的制度与社会后果为它为后续册封、互市或海防安排留下文书与跨政权比较线索。桥接条目只标示可回查的年度问题与材料链；实录有中央选择性，崇祯期尤其需要奏疏、地方、朝鲜及满文材料交叉。

\paragraph{1380（洪武十三年）：胡惟庸密谋刺杀洪武帝，被捕；随后的调查导致大约 15,000 人被处决} 此事把权力、法律或文化关系推进到可辨认的节点。账本的证据限度是编年候选的年序可由官修与后出材料交叉；具体月日、参与者、战果或数字必须回查所列材料，不能将单一叙事当作定论。可资核查的材料为《太祖实录》洪武十三年条；《明史·本纪》及相关列传；诏令、奏疏或会典版本。它们能够说明制度如何表述自身、某种命令或交涉如何被记录，却未必能直接复原家庭内部关系、性别分工、城市日常或海外行动者的意图。事件发生的条件涉及继承、官僚协调或皇权运作的压力使问题进入朝廷程序；其后留下的制度与社会后果为它影响后续人事与政策议程；动机和派系标签不能由单一叙事裁定。译写候选以编年材料定位；实录记载反映朝廷选择，崇祯期须另以奏疏、地方及敌对政权材料交叉。

\subsection{海上联系与外部观察}

\paragraph{1380（洪武十三年）：“黄蜂巢”火箭炮是为明军制造的} 此事把权力、法律或文化关系推进到可辨认的节点。账本的证据限度是编年候选的年序可由官修与后出材料交叉；具体月日、参与者、战果或数字必须回查所列材料，不能将单一叙事当作定论。可资核查的材料为《太祖实录》洪武十三年条；《明史·本纪》及相关列传；诏令、奏疏或会典版本。它们能够说明制度如何表述自身、某种命令或交涉如何被记录，却未必能直接复原家庭内部关系、性别分工、城市日常或海外行动者的意图。事件发生的条件涉及继承、官僚协调或皇权运作的压力使问题进入朝廷程序；其后留下的制度与社会后果为它影响后续人事与政策议程；动机和派系标签不能由单一叙事裁定。译写候选以编年材料定位；实录记载反映朝廷选择，崇祯期须另以奏疏、地方及敌对政权材料交叉。

\paragraph{1381（洪武十四年）十二月：明征服云南：明军占领曲靖} 此事把权力、法律或文化关系推进到可辨认的节点。账本的证据限度是编年候选的年序可由官修与后出材料交叉；具体月日、参与者、战果或数字必须回查所列材料，不能将单一叙事当作定论。可资核查的材料为《太祖实录》洪武十四年条；《明史·兵志》及相关列传；边镇、地方志或对方材料。它们能够说明制度如何表述自身、某种命令或交涉如何被记录，却未必能直接复原家庭内部关系、性别分工、城市日常或海外行动者的意图。事件发生的条件涉及前期边防、地方武装或跨区域资源调动的压力推动了该行动；其后留下的制度与社会后果为它改变了后续调兵、补给或地方秩序的条件；战果和伤亡须分开核对。译写候选以编年材料定位；实录记载反映朝廷选择，崇祯期须另以奏疏、地方及敌对政权材料交叉。

\paragraph{1382（洪武十五年）四月：明征服云南：明军征服云南} 此事把权力、法律或文化关系推进到可辨认的节点。账本的证据限度是编年候选的年序可由官修与后出材料交叉；具体月日、参与者、战果或数字必须回查所列材料，不能将单一叙事当作定论。可资核查的材料为《太祖实录》洪武十五年条；《明史·兵志》及相关列传；边镇、地方志或对方材料。它们能够说明制度如何表述自身、某种命令或交涉如何被记录，却未必能直接复原家庭内部关系、性别分工、城市日常或海外行动者的意图。事件发生的条件涉及前期边防、地方武装或跨区域资源调动的压力推动了该行动；其后留下的制度与社会后果为它改变了后续调兵、补给或地方秩序的条件；战果和伤亡须分开核对。译写候选以编年材料定位；实录记载反映朝廷选择，崇祯期须另以奏疏、地方及敌对政权材料交叉。

\paragraph{1383（洪武十六年）：沿海朝贡、海禁和港口管理的文书构成本年海上联系线索，禁令不能直接推作私人航行停止} 此事把权力、法律或文化关系推进到可辨认的节点。账本的证据限度是来源导向的年度桥接条目：本年材料可定位相关政务线索；具体月日、发文机关、地域和执行结果仍须逐项回查，不能以制度文本或奏报替代实际效果。可资核查的材料为《太祖实录》洪武十六年条；《明史·外国传》；《朝鲜王朝实录》、琉球《历代宝案》或海外档案。它们能够说明制度如何表述自身、某种命令或交涉如何被记录，却未必能直接复原家庭内部关系、性别分工、城市日常或海外行动者的意图。事件发生的条件涉及朝贡名义、边境安全与港口贸易的利益使交涉持续发生；其后留下的制度与社会后果为它为后续册封、互市或海防安排留下文书与跨政权比较线索。桥接条目只标示可回查的年度问题与材料链；实录有中央选择性，崇祯期尤其需要奏疏、地方、朝鲜及满文材料交叉。

\paragraph{1384（洪武十七年）四月：洪武皇帝将政府机构从皇宫迁至南京城墙外} 此事把权力、法律或文化关系推进到可辨认的节点。账本的证据限度是编年候选的年序可由官修与后出材料交叉；具体月日、参与者、战果或数字必须回查所列材料，不能将单一叙事当作定论。可资核查的材料为《太祖实录》洪武十七年条；《明史·本纪》及相关列传；诏令、奏疏或会典版本。它们能够说明制度如何表述自身、某种命令或交涉如何被记录，却未必能直接复原家庭内部关系、性别分工、城市日常或海外行动者的意图。事件发生的条件涉及继承、官僚协调或皇权运作的压力使问题进入朝廷程序；其后留下的制度与社会后果为它影响后续人事与政策议程；动机和派系标签不能由单一叙事裁定。译写候选以编年材料定位；实录记载反映朝廷选择，崇祯期须另以奏疏、地方及敌对政权材料交叉。

\paragraph{1384（洪武十七年）：皇室、功臣与中枢官僚的任免或案件材料构成本年权力重组线索，案狱叙事须避免照录罪名} 此事把权力、法律或文化关系推进到可辨认的节点。账本的证据限度是来源导向的年度桥接条目：本年材料可定位相关政务线索；具体月日、发文机关、地域和执行结果仍须逐项回查，不能以制度文本或奏报替代实际效果。可资核查的材料为《太祖实录》洪武十七年条；《明史·本纪》及相关列传；诏令、奏疏或会典版本。它们能够说明制度如何表述自身、某种命令或交涉如何被记录，却未必能直接复原家庭内部关系、性别分工、城市日常或海外行动者的意图。事件发生的条件涉及继承、官僚协调或皇权运作的压力使问题进入朝廷程序；其后留下的制度与社会后果为它影响后续人事与政策议程；动机和派系标签不能由单一叙事裁定。桥接条目只标示可回查的年度问题与材料链；实录有中央选择性，崇祯期尤其需要奏疏、地方、朝鲜及满文材料交叉。
